\section{One coordinate system for the finite Gaussian comparison}
\label{sec:allarity-coordinate-calibration}

The moment estimate in Section~\ref{sec:allarity-cumulant-comparison}
requires a common space containing both the physical jet coordinates
and the discarded higher-order information. We construct that space
from finite lattice readings. The covariance may be indefinite after
normal ordering; only its symmetry is used in this construction.

\subsection{Recovering polynomial curvature coordinates from differences}

Write $\Delta_a f(x)=f(x+e_a)-f(x)$ on $\mathbb Z^4$. A real
curvature cochain is specified on positively oriented plaquettes and
colour coordinates and extended antisymmetrically to all ordered
pairs of axes. We call it closed when
\[
                \Delta_a f_{bc}+\Delta_b f_{ca}
                                  +\Delta_c f_{ab}=0.
\]
For an integer $d\ge0$, let $K_d$ index a basis of the radial
polynomial connection jets whose curvature degree $j$ satisfies
$j+2\le d$. Its basis element has grade $g_\alpha=j+2$.
Let $t_\alpha$ be its lattice plaquette cochain, obtained by
integrating the polynomial connection along each edge and taking its
discrete curl. These are the polynomial test cochains of the local
section construction. They are closed. Differences of order greater
than $j$ vanish, while order-$j$ differences at zero are the
components of $j!$ times the homogeneous curvature tensor. Edge
integration does not change this highest homogeneous term.

For clarity, the finite sampling argument needs only closedness,
not the divergence equation. Given a closed cochain, its order-$j$
differences at zero form a symmetric derivative tensor. Applying
the closedness equation after any $j-1$ differences gives the
cyclic identity for this tensor. It therefore defines a homogeneous
closed polynomial two-form. For $j=0$ this assertion just says that
the tensor is antisymmetric. The radial primitive is obtained by
contracting this two-form with the position vector and dividing by
$j+2$, as in the local jet construction. Its curl is the given
two-form. Thus these highest-order readings are those of a radial
polynomial connection jet.

The map from the span of the $t_\alpha$ of curvature degree at
least $q$ to all origin differences of orders $q,\ldots,d-2$
is injective. Indeed, take the largest degree with a nonzero
coefficient. The differences of that order kill all smaller
degrees and recover its full curvature tensor. Vanishing of those
readings forces that radial connection jet, hence its basis
coefficients, to be zero. Descend through the finite list of
degrees. The same descent, now subtracting the recovered top-degree
test cochain, shows that the origin readings through order $d-2$
of every closed cochain are matched by a unique linear combination
of the $t_\alpha$.

Choose a linear left inverse of each of these finite injective
sampling maps. For a key $\alpha$ of curvature degree $j$, use the
left inverse involving only differences of orders at least $j$.
This gives a real functional $\ell_\alpha$ which is a finite
linear combination of origin differences, with
\begin{equation}\label{eq:allarity-dual-cochain}
 \ell_\alpha(t_\beta)=\delta_{\alpha\beta},\qquad
 \Pi_df=\sum_{\alpha\in K_d}\ell_\alpha(f)t_\alpha.
\end{equation}
For closed $f$, $\Pi_df$ has exactly its origin differences through
order $d-2$. Compatibility of the different left inverses follows
from uniqueness of the full expansion just proved: restriction to
orders at least $j$ deletes all lower-degree test cochains, and
the left inverse therefore recovers the same remaining coefficients.
The functionals are defined on arbitrary cochains; the matching
assertion is only used for closed ones.

Let $c_\alpha$ be the sum of the absolute values of the coefficients
of $\ell_\alpha$ in its difference representation, and choose
$C_d\ge1$ with $c_\alpha\le C_d$ for every key. This is a finite
constant, selected before any comparison scale or stencil radius.

\subsection{Residual coordinates and their bounds}

Fix a finite set $S$ of addresses, all with lattice distance
$|x|_1\le M$, where $M$ is a nonnegative integer, and a radius $r\ge1$.
For a cochain $f$ define
\begin{align*}
 L_\alpha f&=r^{g_\alpha}\ell_\alpha(f),\qquad \alpha\in K_d,\\
 L_z f&=r^{d+1}(f-\Pi_df)(z),\qquad z\in S.
\end{align*}
The coordinate set is the disjoint union $K_d\sqcup S$. Every
functional $L$ uses only finitely many values, since a finite
difference does so. For every $z\in S$ there is the exact identity
\begin{equation}\label{eq:allarity-coordinate-reconstruction}
 f(z)=\sum_{\alpha\in K_d}r^{-g_\alpha}t_\alpha(z)L_\alpha f
                                      +r^{-(d+1)}L_zf.
\end{equation}
It is simply the definition of $f-\Pi_df$ and holds even when $f$
is not closed.

Put $q=\max(d-1,0)$. Suppose a closed $f$ has the weighted
origin bounds
\[
 r^{j+2}|\Delta_w f(0)|\le P\quad(j=|w|\le d-2)
\]
and the weighted top-difference bounds
\[
 r^{q+2}|\Delta_w f(x)|\le P
             \quad(|w|=q,\ |x|_1\le M).
\]
All plaquette and colour components are included. Assume $P\ge0$.
For physical coordinates, $g_\alpha\le j+2$ for every difference
entering $\ell_\alpha$, so
$|L_\alpha f|\le Pc_\alpha$.

For the residual bound set $h=f-\Pi_df$. Its origin differences
below order $q$ vanish, and its order-$q$ differences equal those
of $f$, since $\Pi_df$ has degree at most $d-2$. If $d=0$ or $1$,
the key set is empty and the last assertion just means $h=f$.
An elementary discrete remainder estimate gives
\[
 |h(z)|\le M^q\sup_{|w|=q,\ |x|_1\le M}|\Delta_w h(x)|.
\]
To prove it, descend from order $q$ to zero. At each lower order
the relevant difference is zero at the origin. Sum its first
differences along a shortest signed coordinate path to $z$.
The path has at most $M$ edges and stays inside the ball. A
negative step uses the forward difference at the new endpoint.
Each descent therefore multiplies the bound by at most $M$.
At $q=0$ no descent is required, with $M^0=1$.
Since $d+1\le q+2$ and $r\ge1$, we obtain
\begin{equation}\label{eq:allarity-coordinate-bounds}
 |L_\alpha f|\le PC_d,\qquad |L_zf|\le PM^q.
\end{equation}
If the weighted origin differences instead have bound $E\ge0$,
the physical coordinates satisfy the sharper bound $EC_d$.

Apply the same argument at both ends of a kernel $K(z,w)$ whose
rows and columns are closed. If its mixed weighted readings of
the preceding types are bounded by $P$, and the origin--origin
ones by $E$, then, for $T\ge\max(C_d,M^q)$,
\begin{equation}\label{eq:allarity-two-reader-bounds}
 |(L_a\otimes L_b)K|\le PT^2,\qquad
 |(L_\alpha\otimes L_\beta)K|\le ET^2.
\end{equation}
Indeed fix one difference at the first endpoint, divide its bound
by its positive weight, and apply
\eqref{eq:allarity-coordinate-bounds} at the second endpoint.
Apply it again at the first. The endpoint operations commute
because both are finite linear combinations of evaluations.
Closedness is preserved under these finite linear combinations.
This justifies both applications, including residual coordinates.

\subsection{The exact free block and the common covariance}

Fix a separated configuration of $n$ insertion points. Let
$F_{i\alpha,j\beta}$ be its free covariance in the radial jet
coordinates. This is a real symmetric matrix: the underlying
free kernel and the change to these coordinates are real. Normal
ordering sets its same-insertion block to zero. Define the reference
cochain kernel
\[
 K^0_{ij}(z,w)=\sum_{\alpha,\beta\in K_d}
 F_{i\alpha,j\beta}\,
 r^{-g_\alpha}t_\alpha(z)\,
 r^{-g_\beta}t_\beta(w).
\]
Its rows and columns are closed. Equation~\eqref{eq:allarity-dual-cochain}
shows that its matrix under the $L$ coordinates is exactly $F$
on the physical block and zero whenever either coordinate is
residual. This is an exact identity, not a small-error statement.

Let $\Gamma_{ij}$ be the actual symmetric cochain kernel and
$\mu_i$ its actual mean. Suppose the means obey the two bounds
of the previous subsection with constants $P,E$, and
$\Gamma_{ij}-K^0_{ij}$ obeys their two-endpoint versions. Define
one enlarged mean and covariance by
\[
 m_{ia}=L_a\mu_i,\qquad
 Q_{ia,jb}=(L_a\otimes L_b)\Gamma_{ij}.
\]
If $T\ge\max(C_d,M^q)$, their entries satisfy
\begin{gather*}
 |m_{ia}|\le PT,\qquad |m_{i\alpha}|\le ET,\\
 |Q_{ia,jb}-F^{\rm ext}_{ia,jb}|\le PT^2,\qquad
 |Q_{i\alpha,j\beta}-F_{i\alpha,j\beta}|\le ET^2,
\end{gather*}
where $F^{\rm ext}$ has zero residual rows and columns. Consequently,
if
\begin{gather}\label{eq:allarity-meridian-purchases}
 B\ge1,\quad PT\le B,\quad PT^2\le B,\quad
 |F_{i\alpha,j\beta}|+PT^2\le B,\\
 ET\le\varepsilon,\qquad ET^2\le\varepsilon,
\end{gather}
all enlarged entries are bounded by $B$, while the physical
block and mean are within $\varepsilon$ of $(F,0)$.
Applying \eqref{eq:allarity-coordinate-reconstruction} at one and
at both endpoints proves that the actual stencil mean and covariance
are the linear and bilinear pullbacks of $(m,Q)$. This gives the
common pullback required for every slot and every subset in the
moment comparison.

Here are the purchases used there. Let $H\to\infty$ be its monitor,
$0<\lambda\le1/4$, and set
\[
 T=P=H^{\lambda/4},\quad E=H^{-3},\quad
 B=H^\lambda,\quad\varepsilon=H^{-2}.
\]
The stencil radius has subpower growth, so $M^q\le T$ eventually;
also $C_d\le T$. Then $PT^2=H^{3\lambda/4}\le B/2$
eventually, while $PT\le PT^2$ since $T\ge1$.
The exponents $-3+\lambda/4$ and $-3+\lambda/2$ are at most $-2$,
which proves the two error inequalities. On a compact separated
configuration set the finitely many free entries are bounded by
continuity of the differentiated kernel; the same bound holds for
the rounded configurations eventually. Thus $|F|\le B/2$
eventually and all of \eqref{eq:allarity-meridian-purchases} holds.

\subsection{The calibrated polynomial price}

The actual polynomial map on a fixed grade space has a coefficient
bound of the form
\[
 |\mathcal H u|_1\le H\,|\operatorname{awake}(u)|_1,
 \qquad |\operatorname{awake}(u)|_1\le2|Tu|_1.
\]
Here $T$ is its tested coefficient map, and $\operatorname{awake}$
is the coefficient vector in the retained common normalization.
For a target label $v$ of grade at most $d$, the calibration equation
is
\[
                D_aTu=\operatorname{weight}_r v,
        \qquad a>0,\quad a^{-1}\le2H,\quad r=H^b\ge1.
\]
$D_a$ multiplies a basis coordinate by $a$ to its connection degree;
$\operatorname{weight}_r$ multiplies it by $r$ to its grade.
Let $q_*$ be the maximum connection degree in the fixed frame, taking
$q_*=0$ if that frame is empty. Coordinatewise inversion gives
\[
 |Tu|_1\le(2H)^{q_*}r^d|v|_1,
 \qquad |\mathcal H u|_1
          \le2^{q_*+1}|v|_1H^{bd+q_*+1}
          \le H^{bd+q_*+2}
\]
eventually. The last threshold depends on the fixed target label,
and can be taken uniformly on a fixed bounded set of labels.
It is uniform in the cutoff and period when the displayed input
bounds are uniform. The finite loss $q_*+2$ is chosen before the
amplification exponent $b$.

Finally the weighted norm of the reconstruction coordinate at
$z\in S$ is exactly
\[
       \sum_{\alpha\in K_d}|t_\alpha(z)|+1.
\]
The powers $r^{-g_\alpha}$ and $r^{-(d+1)}$ cancel their variable
weights. Each $t_\alpha$ is a polynomial of fixed degree at most
$d-2$, so this expression is at most a fixed multiple of
$(1+M)^d+1$. It has subpower growth in $H$. Hence it is eventually
at most $H^\lambda$. If the exact calibrated low-grade identity
and the fixed polynomial degree bound $k$ hold, the weighted tail
estimate of Section~\ref{sec:allarity-cumulant-comparison} applies
with these very coordinates and this coefficient price. Choosing
$b\ge(q_*+2)+3$ leaves its required factor $H^{-2}$.

The construction above proves the coordinate and price conversion.
Its application still requires the actual closed-kernel and mean
difference estimates on the selected atlas, and the exact low-grade
identity for its polynomial map. Those data must refer to the same
readers as the joint Gaussian model. The coordinate conversion
does not obtain them from covariance symmetry alone.
